\documentclass[11pt]{article}
% Uses the canonical latex-docs semantic style in-repository. The local
% lightweight fallback is for independent compilation of this distribution.
\IfFileExists{technical-user-manual.sty}{\usepackage{technical-user-manual}}{\usepackage{../support/aaoa-standalone}}
\usepackage{mathtools}
\setDocTitle{Chapter 5: Analytic Combinatorics}
\setDocSubtitle{Guided reading, derivations and original exercises}
\setDocVersion{1.0}
\setDocStatus{Educational reference / unapproved working material}
\setDocOwner{Algorithm Engineering}
\setDocAuthor{Original companion material}
\setDocSummary{Coverage of primary book Chapter 5 with independently authored examples, proof obligations and lab transfer.}
\providecommand{\Expect}{\mathbb{E}}
\providecommand{\Prob}{\mathbb{P}}
\providecommand{\Var}{\operatorname{Var}}
\providecommand{\bigoh}{\mathcal{O}}
\begin{document}
\makedoctocpage
\section{Reading map and prerequisite check}
\noindent\textbf{Primary source sections:} 5.1 Formal Basis; 5.2 Symbolic Method for Unlabelled Classes; 5.3 Symbolic Method for Labelled Classes; 5.4 Symbolic Method for Parameters; 5.5 Generating Function Coefficient Asymptotics (printed pp. 220--247).\par
\noindent\textbf{Prerequisites:} Completed preceding mathematical-method chapters or equivalent demonstrated mastery.\par
\section{Learning outcomes}
\begin{enumerate}
\item Specify combinatorial class, size and labels without ambiguity.
\item Derive recursive functional equations with a sound bijection.
\item Use singular behavior to infer coefficient growth when hypotheses hold.
\end{enumerate}
\section{Guided technical analysis}
\subsection{From a class to an equation}
For rooted ordered binary trees with $n$ internal nodes, empty-tree choice or root-with-two-subtrees gives $T=1+zT^2$. Solving the quadratic gives a square-root expression; the branch with $T(0)=1$ is required.
\subsection{Asymptotics from the dominant singularity}
The square-root singularity at $z=1/4$ leads to $[z^n]T(z)\sim 4^n/(\sqrt\pi n^{3/2})$. A coefficient-ratio study $C_n\sqrt\pi n^{3/2}/4^n$ should approach one. The ratio does not establish the analytic continuation hypotheses needed by a general transfer theorem.
\subsection{Parameters}
For a statistic $k$, mark objects with $u^k$. Normalized $u$-derivatives at $u=1$ produce mean and factorial moments. Distinguish a uniform random \emph{object} model from a random process that generates objects with unequal probabilities.
\section{Independent problem set}
\begin{enumerate}
\item Prove the tree decomposition is unique and check coefficients for sizes 0--5.
\item Write a two-variable generating function that marks the number of leaves of nonempty binary trees.
\item Construct a case with two equally dominant singularities and explain the parity consequences.
\end{enumerate}
\section{Laboratory transfer}
\textbf{Required laboratory:} LAB-PY03: enumerate ordered binary trees; compare Catalan exact counts and asymptotic predictions.\par
\noindent\textbf{Minimum recorded evidence:} repository commit or source hash; test inputs/seeds; actual operation counts; predicted values; one explicit model limitation.\par
\section{Mastery gate and self-review}
Specification names the atom, labels and size; generating function has correct zeroth coefficient.\par
\begin{enumerate}\item Can I define the random object/input exactly? \item Is my identity exact, asymptotic, or empirical? \item Which base cases and counterexamples did I check? \item What would falsify my performance hypothesis?\end{enumerate}
\section{Suggested notebook structure}
Use five named sections: model and invariants; derivation; deterministic small instances; randomized replication; discrepancy and revised hypothesis. Keep a proof separate from measurement results.

\end{document}
